\section{CNOT-dihedral circuits}\label{app:CNOT-dihedral}

For CNOT-dihedral completeness, the fragment-specific work is narrower than in the other cases.
No encoding/decoding transfer is needed: the source presentation of
\ref{appendix:oldgraphicallanguage5} already lives in a PROP setting with the same scalar
convention as our master signature.

The remaining check is conservative coverage of the source presentation.  Every axiom of
\ref{appendix:oldgraphicallanguage5} omitted from the smaller rule set $\QCdihedral$ must be
derivable from $\QCdihedral$ itself.  The lemmas in this section provide those derivations, so the
imported completeness theorem applies unchanged.

\begin{remark}\label{lem:CNOTdihedral-R4}
\cref{cnot-old-R4} follows by the same rewriting used in
the section for Clifford+T.
\end{remark}

\begin{derivation}\label{lem:CNOTdihedral-R2}
  \cref{cnot-old-R2}
  \begin{gather*}
    \tikzfig{identities/CXC_00}
    \eqeqref{new-R3}\tikzfig{identities/CXC_01}
    \eqeqref{cnot-R1}\tikzfig{identities/CXC_02}\\
    \eqeqref{new-R3}\tikzfig{identities/CXC_03}
    \eqeqref{cnot-R1}\tikzfig{identities/CXC_04}
    \eqeqref{lem:CNOTdihedral-R4}\tikzfig{identities/CXC_05}
    \eqeqref{cnot-R1}\tikzfig{identities/CXC_06}
  \end{gather*}
\end{derivation}

\begin{derivation}\label{lem:CNOTdihedral-src-R3}
  \cref{cnot-old-R3}
  \begin{gather*}
    \tikzfig{cliffordczaxioms/CNOTXCNOT}
    \eqeqref{new-R3}\tikzfig{cliffordczaxioms/CNOTCNOTXX}
    \eqeqref{lem:CNOTdihedral-R4}\tikzfig{cliffordczaxioms/XX}
  \end{gather*}
\end{derivation}

\begin{derivation}\label{lem:CNOTdihedral-src-R5}
  \cref{cnot-old-R5}
  \begin{gather*}
    \tikzfig{cliffordczaxioms/CNOTswCNOTswCNOT}
    \eqeqref{new-R5}\tikzfig{cliffordczaxioms/CnCnNcSwCn}
    \eqeqref{lem:CNOTdihedral-R4}\tikzfig{cliffordczaxioms/NcSwCn}\\
     = \tikzfig{cliffordczaxioms/NcNcSw}
    \eqeqref{lem:CNOTdihedral-R4} \tikzfig{gates/SWAP}
  \end{gather*}
\end{derivation}

\begin{derivation}\label{lem:CNOTdihedral-src-R8}
  \cref{cnot-old-R8}
  \begin{gather*}
    \tikzfig{cliffordczaxioms/R8-old}
    \eqeqref{lem:CNOTdihedral-R4}\tikzfig{cliffordczaxioms/CNOTZCNOT}\\
    \eqeqref{cnot-R7}\tikzfig{cliffordczaxioms/ZZCNOTZCNOT}
    \eqeqref{new-R8} \tikzfig{cliffordczaxioms/CNOTCNOTZZ}
    \eqeqref{lem:CNOTdihedral-R4} \tikzfig{cliffordczaxioms/ZZ}
  \end{gather*}
\end{derivation}

\begin{derivation}\label{lem:CNOTdihedral-src-R11}
  \cref{cnot-old-R11}
  \begin{gather*}
    \tikzfig{cnotdihedral/R11oldL}
    \eqeqref{cnot-R11}\tikzfig{cliffordczaxioms/R11oldmid}
    \eqeqref{cnot-R1} \tikzfig{cnotdihedral/R11oldR}
  \end{gather*}
\end{derivation}

\section{Scalar refinement}\label{app:scalar-refinement}

The scalar-refinement step in \cref{subsec:scalar-refinement} addresses a strict-semantics issue:
global phases are part of the data of the fragment.  Thus two syntactic presentations of the
\emph{same} gate fragment can denote different semantic subPROPs when they generate different
subgroups of global phases.  Before importing a completeness theorem stated with another scalar
convention, we align the scalar subgroups by conservatively extending the syntax with an additional
scalar generator.

The main technical claim proved here is that, under mild assumptions satisfied by all our
unitary fragments, adjoining a new root-of-unity scalar and relating it to an existing one is
a \emph{conservative extension}: it does not introduce any new equalities between circuits
that do not mention the new scalar.  Equivalently, the induced interpretation remains faithful.

%\subsection{Scalars, global phases, and hidden phases}\label{app:scalar-refinement:background}

\begin{definition}\label{def:scalar-in-prop}
Let $\cat{P}$ be a PROP. Its \emph{scalars} are the endomorphisms of the tensor unit,
\(
  S(\cat{P}) \coloneqq \cat{P}(0,0).
\)
Since $0$ is the monoidal unit, both $\circ$ and $\otimes$ restrict to binary operations
\(
  S(\cat{P})\times S(\cat{P})\to S(\cat{P})
\),
and they coincide by interchange.  In particular, $S(\cat{P})$ is a commutative monoid
(and a group whenever $\cat{P}$ is a groupoid).
\end{definition}

\begin{definition}\label{def:global-phase}
Let $ \cat{C}$ be a (strict) symmetric monoidal subcategory of $\cat{FdHilb}$ and assume
$ \cat{C}$ is \emph{endomorphism-only}:
$ \cat{C}(n,m)=\varnothing$ for $n\neq m$.  For $n\in\N$, a unitary
$U\in \cat{C}(n,n)$ is a \emph{global phase on $n$ wires} if it is a scalar multiple of
the identity, namely \(U=\lambda \id_n\) for some
\(\lambda\in\mathrm{U}(1)=\{z\in\C\mid |z|=1\}\).
In words: $U$ acts trivially on all states and only contributes a phase factor $\lambda$.
\end{definition}

\begin{definition}\label{def:visible-scalars}
Let $ \cat{C}\subseteq \cat{FdHilb}$ be a symmetric monoidal subcategory.
The \emph{visible scalar subgroup} of $ \cat{C}$ is
\(
  S( \cat{C}) \coloneqq  \cat{C}(0,0)\subseteq\mathrm{U}(1).
\)
Every element $\lambda\in S( \cat{C})$ yields a global phase on $n$ wires by tensoring:
$\lambda\otimes \id_n = \lambda \id_n\in  \cat{C}(n,n)$.
\end{definition}

\begin{definition}\label{def:hidden-phase}
Let $ \cat{C}\subseteq \cat{FdHilb}$ be an endomorphism-only symmetric monoidal category.
A phase $\lambda\in\mathrm{U}(1)$ is called a \emph{hidden phase of $ \cat{C}$} if there
exists some arity $n\ge 1$ such that $\lambda \id_n\in \cat{C}(n,n)$ while
$\lambda\notin S( \cat{C})$.

Equivalently, $\lambda$ is hidden if it occurs as a global phase at some positive wire count,
but cannot be produced as a scalar morphism $0\to 0$ inside~$ \cat{C}$.
\end{definition}

\begin{example}\label{ex:hidden-phase}
Consider the subgroup $ \cat{C}\subseteq \cat{Qubit}$ generated (as a symmetric monoidal
subcategory) by a single 1-qubit gate $R_X(2\pi)$, but \emph{with no scalar generators}.
Semantically, $R_X(2\pi)=-\id_1$, hence $-\id_1\in \cat{C}(1,1)$.  However,
$S( \cat{C})=\{1\}$ because there are no nontrivial $0\to 0$ morphisms available.
Thus $\lambda=-1$ is a hidden phase in the sense of \cref{def:hidden-phase}.

If we were to adjoin a new scalar generator $\omega$ interpreted as $-1$, then $\omega\otimes\id_1$
and the existing circuit $R_X(2\pi)$ would have the same semantics; without additional axioms
linking them, this would typically destroy faithfulness of the interpretation.  The ``no hidden phases''
assumption in \cref{lem:scalar-refinement} precisely rules out this situation.
\end{example}

%\subsection{A categorical fact: scalars are central}\label{app:scalar-refinement:centrality}

The conservativity argument relies on the observation that scalars commute with all endomorphisms,
so new scalar generators can always be ``pulled out'' as a single global phase factor.

\begin{lemma}\label{lem:scalar-centrality}
Let $ \cat{C}$ be a strict symmetric monoidal category.  For any scalar $a:0\to 0$ and any
endomorphism $f:n\to n$ one has
\begin{equation}\label{eq:scalar-centrality-app}
  (a\otimes \id_n)\circ f = f\circ (a\otimes \id_n).
\end{equation}
For scalars $a,b:0\to 0$, composition and tensor coincide:
\(
  a\circ b = a\otimes b.
\)
\end{lemma}
\begin{proof}
Since the tensor is bifunctorial and the category is strict,
\((a\otimes \id_n)\circ f
= (a\circ \id_0)\otimes(\id_n\circ f)
= (\id_0\circ a)\otimes(f\circ \id_n)
= f\circ(a\otimes \id_n)\), using the interchange law.
For scalars, the same computation with $n=0$ yields $a \circ b = a \otimes b$.
\end{proof}

%\subsection{The refinement construction and its intended semantics}\label{app:scalar-refinement:construction}

The scalar-refinement step is an extension of a presentation by adjoining a new scalar generator of
higher order.

\begin{definition}\label{def:mu_m}
For $m\ge 1$ write
\(
  \mu_m \coloneqq \{e^{2\pi i k/m}\mid 0\le k<m\}\subseteq \mathrm{U}(1)
\)
for the cyclic group of $m$-th roots of unity.
\end{definition}

\begin{definition}\label{def:scalar-refined-presentation}
Let $\cat{P}^{\mathrm{src}}/\mathcal{R}^{\mathrm{src}}$ be a presented PROP and let
$s:0\to 0$ be a chosen scalar in the \emph{src} signature.
Fix integers $m\ge 1$ and $\ell\ge 1$.
Define $\cat{P}^{\mathrm{src},\sharp}/\mathcal{R}^{\mathrm{src},\sharp}$ to be the
presentation obtained by adjoining a new scalar generator
\(
  \omega:0\to 0
\)
and adding the two relations
\begin{equation}\label{eq:refinement-relations}
  \omega^{m\ell}=\id_0,
  \qquad
  \omega^{r}=s,
\end{equation}
where $\omega^k$ denotes the $k$-fold composite of $\omega$ (equivalently, $k$-fold tensor),
and $r$ is an integer chosen so that semantically $\omega^r$ matches $s$ as in
\eqref{eq:choose-r}.
\end{definition}

\begin{remark}\label{rem:refinement-intuition}
The relations \eqref{eq:refinement-relations} assert that $\omega$ has order dividing $m\ell$,
and that the \emph{src} scalar $s$ is the $r$-th power of $\omega$.  Concretely, this is the
standard ``adjoin a root'' construction at the level of scalars.  The point of
\cref{lem:scalar-refinement} is that, in a PROP setting, adjoining such a root is
conservative provided there are no hidden phases.
\end{remark}

%\subsection{A normal form: extracting the new scalar}\label{app:scalar-refinement:normal-form}

The centrality of scalars yields a simple normal form in the refined syntax: all occurrences of
$\omega$ can be extracted into a single scalar factor.

\begin{lemma}\label{lem:scalar-extraction}
In the refined PROP $\cat{P}^{\mathrm{src},\sharp}/\mathcal{R}^{\mathrm{src},\sharp}$,
every morphism $C:n\to n$ is derivably equal to one of the form
\begin{equation}\label{eq:scalar-normal-form}
  C = \omega^{a}\otimes C',
\end{equation}
where $a\in\{0,1,\dots,m\ell-1\}$ and $C'$ is a morphism built using only the
\emph{src} generators (i.e.\ $C'$ contains no occurrence of $\omega$).
\end{lemma}
\begin{proof}
Work in the free PROP on the refined signature, modulo PROP coherence and the relations
$\mathcal{R}^{\mathrm{src}}$ and \eqref{eq:refinement-relations}.  Because $\omega$ is a scalar,
every occurrence of $\omega$ in a diagram is a $0\to 0$ component tensored somewhere with an
endomorphism.  By \cref{lem:scalar-centrality}, we can slide each such component past all other
boxes and past all compositions/tensors so as to collect all $\omega$-factors together as a single
power $\omega^{a}$ tensored with a diagram $C'$ that contains no $\omega$.
Finally, use $\omega^{m\ell}=\id_0$ to reduce $a$ modulo $m\ell$ into the stated range.
\end{proof}

%\subsection{Conservativity of scalar refinement}\label{app:scalar-refinement:main}

The main lemma used in \cref{subsec:scalar-refinement} is the resulting conservativity statement.

\begin{lemma}\label{lem:scalar-refinement}
Let $\cat{P}^{\mathrm{src}}/\mathcal{R}^{\mathrm{src}}$ be a graphical language equipped with a
\emph{faithful} (i.e.\ complete) strict symmetric monoidal interpretation
\(
  \interp{\cdot}_{\mathrm{src}} :
  \cat{P}^{\mathrm{src}}/\mathcal{R}^{\mathrm{src}}
  \longrightarrow
   \cat{C}^{\mathrm{src}}
  \subseteq \cat{FdHilb}.
\)
Assume the following semantic properties.
\begin{enumerate}
\item\label{it:scalar-cyclic-app}
(finite cyclic visible scalars)
The scalar group $S( \cat{C}^{\mathrm{src}})= \cat{C}^{\mathrm{src}}(0,0)$ is finite cyclic of order $m$
and generated by $\interp{s}_{\mathrm{src}}$ for some chosen src scalar $s:0\to 0$.
Equivalently, $S( \cat{C}^{\mathrm{src}})\cong \mu_m$.

\item\label{it:no-hidden-phases-app}
(no hidden phases)
$ \cat{C}^{\mathrm{src}}$ has no hidden phases: for all $n\in\N$ and all $\lambda\in\mathrm{U}(1)$,
if $\lambda \id_n\in \cat{C}^{\mathrm{src}}(n,n)$ then $\lambda\in S( \cat{C}^{\mathrm{src}})$
(cf.\ \cref{def:hidden-phase}).

\item\label{it:invertible-app}
(invertibility)
Every morphism in $ \cat{C}^{\mathrm{src}}$ is invertible, i.e.\ each hom-set
$ \cat{C}^{\mathrm{src}}(n,n)$ is a group under composition.\footnote{%
In all applications in this paper, $ \cat{C}^{\mathrm{src}}$ consists of unitaries, hence this holds automatically.}
\end{enumerate}

Fix $\ell\ge 1$ and choose a primitive root of unity $\zeta\in\mathrm{U}(1)$ of order $m\ell$.
Choose an integer $r$ such that
\begin{equation}\label{eq:choose-r}
  \zeta^{r}=\interp{s}_{\mathrm{src}}.
\end{equation}
Form the scalar-refined presentation
$\cat{P}^{\mathrm{src},\sharp}/\mathcal{R}^{\mathrm{src},\sharp}$ as in
\cref{def:scalar-refined-presentation}.

Let $ \cat{C}^{\sharp}\subseteq\cat{FdHilb}$ be the subPROP generated by $ \cat{C}^{\mathrm{src}}$
together with the scalar $\zeta:0\to 0$ (viewed as the unitary map $\C\to\C$, $1\mapsto\zeta$).

Then the interpretation functor extends uniquely to a strict symmetric monoidal functor
\(
  \interp{\cdot}^{\sharp} :
  \cat{P}^{\mathrm{src},\sharp}/\mathcal{R}^{\mathrm{src},\sharp}
  \longrightarrow
   \cat{C}^{\sharp}
\)
sending $\omega$ to $\zeta$, and $\interp{\cdot}^{\sharp}$ is faithful.

In particular, the refined language is complete for $ \cat{C}^{\sharp}$ and
$S( \cat{C}^{\sharp})\cong \mu_{m\ell}$.
\end{lemma}

\begin{proof}
(1) Soundness and existence of the refined interpretation.
By construction, $\interp{\cdot}^{\sharp}$ agrees with $\interp{\cdot}_{\mathrm{src}}$ on all src generators.
We interpret the new generator $\omega$ by the scalar $\zeta$.
The additional relations \eqref{eq:refinement-relations} hold semantically because
$\zeta^{m\ell}=1$ and $\zeta^{r}=\interp{s}_{\mathrm{src}}$ by \eqref{eq:choose-r}.
Therefore the universal property of presented PROPs yields a unique induced functor
$\interp{\cdot}^{\sharp}$ into $ \cat{C}^{\sharp}$.

(2) Scalar extraction in the refined syntax.
By \cref{lem:scalar-extraction}, every refined circuit $C:n\to n$ is derivably equal to
$\omega^{a}\otimes C'$ with $0\le a<m\ell$ and $C'$ an src circuit.

(3) Faithfulness.
Let $C_1,C_2:n\to n$ be refined circuits such that
$\interp{C_1}^{\sharp}=\interp{C_2}^{\sharp}$.
Write both in scalar-extraction normal form:
\(
  C_i = \omega^{a_i}\otimes C_i'
  \qquad (0\le a_i<m\ell),
\)
with $C_i'$ containing no $\omega$.
Let $U_i\coloneqq \interp{C_i'}_{\mathrm{src}}\in \cat{C}^{\mathrm{src}}(n,n)$.
The semantic equality becomes
\begin{equation}\label{eq:semantic-equality-expanded}
  \zeta^{a_1}   U_1 = \zeta^{a_2}  U_2.
\end{equation}
By invertibility \ref{it:invertible-app}, we can right-multiply by $U_2^{-1}$ to obtain
\begin{equation}\label{eq:phase-in-src-app}
  \zeta^{a_2-a_1} \id_n = U_1\circ U_2^{-1}
  \in  \cat{C}^{\mathrm{src}}(n,n).
\end{equation}
Thus $\zeta^{a_2-a_1}$ appears as a global phase in $ \cat{C}^{\mathrm{src}}$ on $n$ wires.
By the no-hidden-phases assumption \ref{it:no-hidden-phases-app}, this implies
\(
  \zeta^{a_2-a_1}\in S( \cat{C}^{\mathrm{src}}).
\)
By \ref{it:scalar-cyclic-app}, there exists $t\in\{0,\dots,m-1\}$ such that
\begin{equation}\label{eq:phase-as-src-scalar}
  \zeta^{a_2-a_1} = \interp{s}_{\mathrm{src}}^{ t}.
\end{equation}
Combining \eqref{eq:semantic-equality-expanded} with \eqref{eq:phase-as-src-scalar} gives
\(
  \interp{C_1'}_{\mathrm{src}}
  =
  \interp{s^{t}\otimes C_2'}_{\mathrm{src}}.
\)
Since $\interp{\cdot}_{\mathrm{src}}$ is faithful by hypothesis,
\(
  \cat{P}^{\mathrm{src}}/\mathcal{R}^{\mathrm{src}}
  \vdash
  C_1' = s^{t}\otimes C_2'.
\)
The refined presentation contains all src relations, so the same derivation is valid in
$\cat{P}^{\mathrm{src},\sharp}/\mathcal{R}^{\mathrm{src},\sharp}$.
Now use the refinement relation $\omega^{r}=s$ to rewrite $s^{t}=\omega^{rt}$, yielding
\(
  C_1
  =
  \omega^{a_1}\otimes C_1'
  =
  \omega^{a_1}\otimes \omega^{rt}\otimes C_2'
  =
  \omega^{a_1+rt}\otimes C_2'.
\)

It remains to compare the exponents.  From \eqref{eq:phase-as-src-scalar} and \eqref{eq:choose-r} we have
$\zeta^{a_2-a_1}=\zeta^{rt}$, so $\zeta^{a_2-a_1-rt}=1$.  Since $\zeta$ has exact order $m\ell$, this implies
\(
  a_2-a_1-rt \equiv 0 \pmod{m\ell}.
\)
Using $\omega^{m\ell}=\id_0$ in the refined syntax, we conclude $\omega^{a_1+rt}=\omega^{a_2}$, hence
\(
  C_1=\omega^{a_2}\otimes C_2' = C_2.
\)
Therefore $\interp{\cdot}^{\sharp}$ is faithful.

(4) Scalars of $ \cat{C}^{\sharp}$.
By definition, $S( \cat{C}^{\sharp})$ is generated by $S( \cat{C}^{\mathrm{src}})=\langle \interp{s}_{\mathrm{src}}\rangle$
together with $\zeta$.  Since $\interp{s}_{\mathrm{src}}=\zeta^r\in\langle \zeta\rangle$, we obtain
$S( \cat{C}^{\sharp})=\langle \zeta\rangle\cong \mu_{m\ell}$.
\end{proof}

\begin{corollary}\label{cor:scalar-refinement-conservative}
With the hypotheses of \cref{lem:scalar-refinement}, let $C_1,C_2$ be \emph{src} circuits
(i.e.\ containing no occurrence of $\omega$).  If
\(
  \cat{P}^{\mathrm{src},\sharp}/\mathcal{R}^{\mathrm{src},\sharp}\vdash C_1=C_2,
\)
then already
\(
  \cat{P}^{\mathrm{src}}/\mathcal{R}^{\mathrm{src}}\vdash C_1=C_2.
\)
\end{corollary}
\begin{proof}
If the refined theory derives $C_1=C_2$, then in particular $\interp{C_1}_{\mathrm{src}}=\interp{C_2}_{\mathrm{src}}$
(because the refined interpretation restricts to the src one on src syntax).  Faithfulness of
$\interp{\cdot}_{\mathrm{src}}$ implies the src derivability.
\end{proof}

%\subsection{Instances used in this paper}\label{app:scalar-refinement:instances}

The two scalar refinements applied in the main text are the following.

\begin{itemize}
\item \textbf{Qutrit Clifford.}
The imported presentation of \cite{qutritclifford} uses an order-$6$ scalar
$\tikzfig{gates/minw}$.  Our strict qutrit convention uses an order-$12$ scalar subgroup
(\cref{sec:circuits-relations}).  We refine the source syntax by adjoining an order-$12$ scalar
generator $\tikzfig{gates/wb}$ and adding the relations
\(\tikzfig{gates/wb}^{12}=\id_0\) and
\(\tikzfig{gates/wb}^{10}=\tikzfig{gates/minw}\).
This is the $\mu_6\to\mu_{12}$ refinement referenced in \cref{subsec:scalar-refinement}.

\item \textbf{Clifford+CS.}
The imported presentation of \cite{SelingerCliffordPlusCS} uses an order-$4$ scalar
$\tikzfig{gates/i}$, whereas our strict qubit convention uses an order-$8$ scalar subgroup.
We refine the source syntax by adjoining an order-$8$ scalar generator $\tikzfig{gates/wprime}$
and adding the relations \(\tikzfig{gates/w}^{8}=\id_0\) and
\(\tikzfig{gates/w}^{2}=\tikzfig{gates/i}\).
This is the $\mu_4\to\mu_8$ refinement referenced in \cref{subsec:scalar-refinement}.
\end{itemize}

In both cases, the hypotheses of \cref{lem:scalar-refinement} are satisfied for the unitary
semantics considered in this paper: invertibility is automatic, visible scalars are finite cyclic by
construction, and the imported normal forms expose every scalar multiple of an identity as one of
those visible scalars, so no hidden phases appear relative to the chosen scalar conventions.
Therefore, scalar refinement preserves faithfulness, and the refined imported presentations may be
treated as complete presentations for the strict semantics used throughout.


%\section{Minimality interpretations for all fragments}


\begin{figure}[ht]
    \begin{threeparttable}
    \begin{adjustbox}{width=\textwidth,center}
    \renewcommand{\arraystretch}{1.2}
    \begin{tabular}{|c|c|c|c|c|c|c|c|}
    \hline
    \textbf{Interpretation} & \textbf{Cliff} & \textbf{Real Cliff} & \textbf{Qutrit Cliff} & \textbf{Cliff+T} & \textbf{Cliff+CS} & \textbf{CNOT-dih} &  \textbf{Arity} \\
    \hline
    $?\gw$ & \eqref{w8} & — & — & \eqref{t-w8}  & \eqref{cs-w8} & \eqref{cnot-R10} & 0\\
    $?\gminus$ & — & \eqref{realminus2} & — & —  & — & — & 0\\
    $?\gww$ & — & — & \eqref{qt-omegapow12} & —  & — & — & 0\\
    \hline
    $?\gH$ & \eqref{H2} & \eqref{realH2} & \eqref{qt-hpow4} & \eqref{t-H2}  & \eqref{cs-H2} & — & 1 \\
    $?\gS$ & \eqref{S4} & — & \eqref{qt-spow3} & — & \eqref{cs-S4} & — & 1 \\
    $?\gZ$ & — & \eqref{realZ2} & — & — & — & — & 1 \\
    $?\gT$ & — & — & — & \eqref{t-Tpow8} & — & \eqref{cnot-R7} & 1 \\
    $?\gX$ & — & — & — & — & — & \eqref{cnot-R1} & 1 \\
    $\#\{\gH\}_{[2]}$ & \eqref{E} & — & \eqref{qt-shpow3} & \eqref{t-E} & \eqref{cs-E} & — & 1 \\
    $\#\{\gw\}_{[2]}$ & — & — & — & — & — & \eqref{cnot-R11} & 1 \\
    $\#\{\gminus\}_{[2]}$ & — & \eqref{realF} & — & — & — & — & 1 \\
    $\#\{\gH, \gw\}_{[2]}$ & — & — & — & \eqref{t-TX} & — & — & 1 \\
    $\#\{\gX\}_{[2]}$ & — & — & — & — & — & \eqref{new-R3} & 1 \\
    $[\gS:=\gSX]_{\sim}$ & — & — & \eqref{qt-ssprime}\tnote{a} & — & — & — & 1 \\
    \hline
    $?\gCNOT$ & \eqref{Cs} & \eqref{realCX2} & \eqref{qt-cnotremove} & \eqref{t-Ct} & — & \eqref{cnot-R12} & 2 \\
    $?\gCS$ & — & — & — & — & \eqref{cs-C} & — & 2 \\
    $\#\{\gSWAP\}_{[2]}$ & \eqref{B} & \eqref{realB} & \eqref{qt-swapdecomp} & \eqref{t-B} & \eqref{cs-B} & \eqref{new-R5} & 2 \\
    $\#\{\gCNOT,\gSWAP\}_{[2]}$ & \eqref{CZ} & —  & — & \eqref{t-CZ} & — & — & 2 \\
    $\#\{\gS\}_{[3]}$ & — & — & \eqref{qt-czdecomp} & — & — & — & 2 \\
    $\#\{\gZ\}_{[2]}$ & — & \eqref{realwC} & — & — & — & — & 2 \\
    $[\gZ:=\gI]_{\sim}$ & — & \eqref{realCF} & — & — & — & — & 2 \\
    $[\gH:=\gI]_{\sim}$ & — & \eqref{realXC} & — & — & — & — & 2 \\
    $\#\{\gCNOT\}_{[2]}$ & — & — & \eqref{qt-kcnot} & — & — & — & 2 \\
    $\#\{\gS,\gH\}_{[2]}$ & — & — & — & — & \eqref{cs-XCS} & — & 2 \\
    $[\gCS:=\gCSz]_{\sim}$ & — & — & — & — & \eqref{cs-CSrev} & — & 2 \\
    $[\gCS:=\gCSzz]_{\sim}$ & — & — & — & — & \eqref{cs-SHCHC} & — & 2 \\
    $\#\{\gT,\gw,\gw\}_{[8]}$ & — & — & — & — & — & \eqref{new-R8} & 2 \\
    \hline
    $\arg\det_{2}$ & \eqref{Inew} & \eqref{realInew} & — & — & \eqref{cs-I} & — & 3 \\
    $\#\{\gCNOT,\gSWAP\}_{[2]}$ & — & —  & — & — & — & \eqref{cnot-R6} & 3 \\
    $\#\{\gT,\gw,\gw\}_{[4]}$ & — & — & — & — & — & \eqref{cnot-R9} & 3 \\
    \hline
    $\arg\det_{3}$ & — & — & — & — & — & \eqref{cnot-R13} & 4 \\
    \hline
    \end{tabular}
    \end{adjustbox}
    \begin{tablenotes}
    \footnotesize
    \item[a] See \cref{qt-prop:SSpnecessity}.
    \end{tablenotes}
    \end{threeparttable}
  \caption{Summary of minimality interpretations across all fragments}
  \label{fig:minimality-summary}
\end{figure}
